\documentclass[11pt]{article}
\usepackage[margin=1in]{geometry}
\usepackage{amsmath,amssymb,amsthm}
\usepackage[T1]{fontenc}
\usepackage[colorlinks=true,linkcolor=blue,citecolor=blue,urlcolor=blue]{hyperref}
\newtheorem{theorem}{Theorem}
\newtheorem{lemma}[theorem]{Lemma}
\newtheorem{proposition}[theorem]{Proposition}
\newtheorem{remark}[theorem]{Remark}
\newcommand{\F}{\mathbb F}
\newcommand{\E}{\mathcal E}
\newcommand{\half}{\tfrac12}
\title{Sparse halves for Clebsch blow-ups:\\six exceptional classes and a star extension}
\author{\texttt{kasaharamihano-spec}}
\date{24 October 2026\\[3pt]\small Research draft; not peer reviewed}
\begin{document}
\maketitle
\begin{abstract}
Let $H$ be the Clebsch graph, and replace each vertex $v$ by an independent
set of size $w_v$, with complete bipartite graphs on the edges of $H$.
We prove the sparse-half bound $N^2/50$ whenever at most six class sizes
differ from a common value. The proof reduces the exceptional positions
to six graph shapes using the five-generator representation, then treats
their arbitrary capacities by matching cuts, weighted pentagon
constructions, and Razborov's independence-ratio theorem.
An explicit square identity gives a pentagon half avoiding a prescribed
edge and supports the final two cyclic cases.
We also treat seven arbitrary capacities on a vertex, its five neighbours,
and one further vertex, with the other nine capacities equal. A direct
lift preserves the constructions for independently unequal leaf capacities.
Finally, an elementary matching criterion permits all sixteen positive
capacities to vary when their maximum-to-minimum ratio is at most
$(24+5\sqrt{21})/17$.
This does not settle arbitrary Clebsch blow-ups or the general
sparse-half conjecture. The priority of this restricted consequence
has not been established.
\end{abstract}

\section{Statements and scope}
Use the model
\[
 V(H)=\F_2^4,\qquad
 D=\{1,2,4,8,15\},\qquad uv\in E(H)\ \Longleftrightarrow\ u+v\in D ,
\]
where integers encode four-bit vectors and addition is bitwise XOR.
The only nonempty subset of $D$ with sum zero is $D$ itself. In particular
$H$ is triangle-free. For capacities $w_v\geq0$, put $N=\sum_v w_v$.
A fractional half is a vector $x$ with
\[
 0\leq x_v\leq w_v,\qquad \sum_vx_v=N/2,\qquad
 \E(x)=\sum_{uv\in E(H)}x_ux_v .
\]
For integer capacities this is the natural relaxation of selecting vertices
from the corresponding blow-up.

\begin{theorem}\label{thm:main}
If there is $c\geq0$ such that at most six capacities $w_v$ differ from $c$,
the Clebsch blow-up admits a fractional half of cost at most $N^2/50$.
For integer capacities it has a set of $\lfloor N/2\rfloor$ vertices
inducing at most $N^2/50$ edges.
\end{theorem}

We prove the following family statements first, and give a short
structural reduction to them at the end. The numbered parts below
are referenced in the individual proofs.

\begin{proposition}\label{prop:families}
There is a fractional half of cost at most $N^2/50$ in each of the following
families of nonnegative real capacities.
\begin{enumerate}
\item There is $c\geq0$ such that at most five capacities differ from $c$.
\item The capacities on $\{z\}\cup N_H(z)$ are arbitrary, and all other
ten capacities equal $c\geq0$.
\item The capacities on an induced $C_5=P$ and a vertex $z$ nonadjacent
to every vertex of $P$ are arbitrary, and all other ten capacities equal
$c\geq0$.
\item The capacities on an automorphic image of
$\{1,2,4,6,8,9\}$ are arbitrary, and all other ten capacities equal
$c\geq0$. These six vertices induce a double star with adjacent centres
$6,9$, left leaves $2,4$, and right leaves $1,8$.
\item The capacities on $\{0,1,2,4,7,8\}$ are arbitrary, with the
remaining ten equal. This is a star with one leaf edge extended.
\item The capacities on $\{0,1,3,4,7,8\}$ are arbitrary, with the
remaining ten equal. This is the pentagon $(0,1,3,7,8)$ with a
pendant vertex $4$ attached to $0$.
\item The capacities on $\{0,1,2,3,7,8\}$ are arbitrary, with the
remaining ten equal. This is a theta graph with two paths of length two
and one path of length three between $0$ and $3$.
\end{enumerate}
For integer capacities, the blow-up has a set of $\lfloor N/2\rfloor$
vertices inducing at most $N^2/50$ edges.
\end{proposition}

Our external inputs are explicitly identified below. In particular,
the weighted $C_5$ construction is already present in
Norin--Yepremyan~\cite[proof of Theorem 3.1, $d=2$]{NY}.
The independence-ratio input is Razborov~\cite[Corollary 3.7]{R}.
No numerical experiment is used as a proof of a continuous statement.

\section{Preliminaries}
\begin{lemma}[Matching cuts]\label{lem:cuts}
For each $d\in D$ there is a linear functional $\ell_d:\F_2^4\to\F_2$
that takes value $0$ on $d$ and value $1$ on $D\setminus\{d\}$.
Each fibre induces four disjoint edges, all of difference $d$.
\end{lemma}
\begin{proof}
Prescribing the values is consistent with the unique relation
$\sum_{s\in D}s=0$, since four prescribed values are $1$.
The functional is nonzero. Each fibre has eight vertices; the only
generator preserving it is $d$.
\end{proof}

Call the ten fibres in Lemma~\ref{lem:cuts} matching sides.
The case $N=0$ is immediate; henceforth assume $N>0$.
If a side has total capacity $b\geq N/2$ and full cost $Q\leq b^2/16$,
scaling all its capacities by $N/(2b)$ gives a half with cost at most
$N^2/64$. We will repeatedly use this sufficient condition.

\begin{lemma}[Capacitated matching]\label{lem:matching}
For disjoint edges with capacities $(L_j,S_j)$, $L_j\geq S_j$,
a minimum-cost selection of prescribed total mass is obtained by first
using the larger endpoints at cost zero, then filling smaller endpoints
in increasing order of $L_j$, at cost $L_j$ per unit.
Moreover, at fixed total selection mass the minimum cost is nonincreasing
when any capacity increases.
\end{lemma}
\begin{proof}
On a single edge, fixing its selected mass $t$ gives a concave quadratic
$x(t-x)$. An endpoint of its feasible interval minimizes this expression.
The resulting minimum is zero for $0\leq t\leq L_j$ and
$L_j(t-L_j)$ for $L_j\leq t\leq L_j+S_j$.
Allocation over the edges is then a linear filling problem.
The last assertion follows directly by inclusion of feasible sets:
the objective does not depend on the capacities.
\end{proof}

\begin{lemma}[Known independence-ratio input]\label{lem:alpha}
If an independent set in a triangle-free weighted core has total capacity
at least $2N/5$, the weighted core admits a fractional half of cost at most $N^2/50$.
\end{lemma}
\begin{proof}
For integer capacities, double every capacity. The resulting graph has
$2N$ vertices and independence number at least $4N/5$.
Razborov's theorem~\cite[Corollary 3.7]{R}, including its weighted-half
formulation, gives a half of cost at most $(2N)^2/50$.
Dividing the selected class masses by two proves the claim.
Rational capacities follow by clearing denominators.
For real capacities, normalize $N=1$ and approximate within the rational
polytope defined by nonnegativity and the specified independent-set
inequality. Compactness of the selected vectors and continuity of the
cost give the conclusion. The case $N=0$ is immediate.
\end{proof}

\begin{lemma}[Rounding]\label{lem:rounding}
For integer capacities, a fractional half of cost at most $N^2/50$
implies the integer assertion in Theorem~\ref{thm:main}.
\end{lemma}
\begin{proof}
The two-coordinate mechanism is classical; see, for example,
Brandt~\cite[Proposition 1]{B}. We include the capacitated argument.
If $N$ is odd, delete mass $1/2$; all cost coefficients are nonnegative,
so cost cannot increase. The target mass is now $\lfloor N/2\rfloor$.
Given two coordinates strictly between zero and their integer capacities,
vary them as $x_i+t,x_j-t$. The cost is concave in $t$, its quadratic
coefficient being $-a_{ij}\leq0$. Some feasible endpoint does not increase
cost and moves at least one further coordinate to a capacity boundary.
Repeat until at most one coordinate is internal.
Its integrality follows from the integer total mass.
\end{proof}

\begin{lemma}[Known pentagon construction]\label{lem:c5}
For nonnegative capacities $b_0,\ldots,b_4$ on a cycle, with sum $T$,
there is a half of cost at most $T^2/50$ supported in
$\{i,i+1,i+3\}$ for some $i$ (indices are modulo five).
\end{lemma}
\begin{proof}
This is the $d=2$ construction in~\cite[proof of Theorem 3.1]{NY},
where the authors observe that the minimum-degree hypothesis is not
used in this branch. We recall its short proof in cycle order.
If a nonadjacent pair has mass at least $T/2$, use it at zero cost.
Every such pair lies in a triple of the required form.
Otherwise fill $b_i,b_{i+3}$ and take
$T/2-b_i-b_{i+3}$ from $i+1$. All five choices are feasible because
both the filled nonadjacent pair and the omitted nonadjacent pair
have mass less than $T/2$. Their costs $q_i$ satisfy
\[
 \sum_iq_i=\frac{T^2}{2}
   -\frac12\sum_i(b_i+b_{i+3})^2
 \leq \frac{T^2}{10},
\]
by Cauchy--Schwarz.
\end{proof}

\section{At most five exceptional classes}
Let $X=\{v:w_v\ne c\}$. If a generator $d$ is absent from the differences
of edges internal to $X$, every $d$-edge has an endpoint of capacity $c$.
Take the heavier matching side, of capacity $b\geq N/2$.
Its full cost is $Q=cb-4c^2$, so
\[
 b^2-16Q=(b-8c)^2\geq0.
\]
The scaling observation proves the required bound, including $c=0$.

If $|X|\leq5$ and all five differences occur, choose one edge of each
difference within $X$. This gives five edges on at most five vertices,
so it contains a cycle. The distinct differences of that cycle sum to
zero and hence comprise all of $D$. Thus $X$ is exactly an induced
pentagon. It remains to treat arbitrary capacities $p_0,\ldots,p_4$ on
a pentagon and common capacity $c$ on the other eleven vertices.
Put $W=\sum p_i$ and $N=W+11c$.

Along an induced pentagon each generator occurs once.
For its edge $\{i,i+1\}$, the corresponding matching cut splits the
pentagon into $\{i,i+1,i+3\}$ and $\{i+2,i+4\}$.
If some $p_0\leq c$, one of the nonadjacent pairs
$\{1,3\},\{2,4\}$ has capacity $s\geq(W-c)/2$.
The matching side meeting the pentagon in that pair has
\[
 b=s+6c\geq N/2,\quad Q=cs+2c^2,\quad
 b^2-16Q=(s-2c)^2 .
\]
It therefore gives cost at most $N^2/64$.

Otherwise put $b_i=p_i-c\geq0$ and $T=\sum b_i$, so $N=T+16c$.
Use Lemma~\ref{lem:c5}, and take the matching side containing its
support triple. Select $c$ at each of the eight vertices of this side
and add the half from the $b_i$. The mass is $N/2$ and the cost is at most
\[
 4c^2+\frac{cT}{2}+\frac{T^2}{50}\leq\frac{N^2}{50},
\]
since the difference is $28c^2/25+7cT/50$.
This proves part (1).

\section{A vertex and its five neighbours}
Translate the centre to $0$, give it capacity $d$, give its five neighbours
capacities $a_1,\ldots,a_5$, and give the remaining ten vertices capacity $c$.
Write $S=\sum a_i$, $a=S/5$, and $N=d+S+10c$.
The matching side $A_i$ containing the centre and leaf $i$ consists of
\[
 (d,a_i),(c,c),(c,c),(c,c).
\]
Its complement $B_i$ consists of $(a_j,c)$, $j\ne i$.
Two independent-set witnesses have capacities $S$ and $d+4c$.

We first give the symmetric calculation needed in the proof.
\begin{lemma}[Symmetric matching-side bound]\label{lem:star}
Assume $a>0$, $N=d+5a+10c$,
\[
 5a<2N/5,\qquad d+4c<2N/5,\qquad d>3a-2c .
\]
The matching with capacities $(d,a)$ and three copies of $(c,c)$
contains a half of cost at most $N^2/50$.
\end{lemma}
\begin{proof}
Set $x=d/a$, $t=c/a$, $n=N/a=x+5+10t$.
The assumptions imply
\[
 0\leq x<10/3,\qquad x>15/2-10t,\qquad x>3-2t. \tag{R}
\]
The last inequality makes the side heavier than $N/2$.
Let $H$ be its exact minimum cost divided by $a^2$, and
$\Delta=n^2-50H$.

If $t\geq\max(x,1)$, fill both star endpoints before paying for background
small endpoints. The remaining background demand is $(3+4t-x)/2>0$,
and Lemma~\ref{lem:matching} gives
\[
 H=x+t(3+4t-x)/2,\quad
 \Delta=x^2-40x+25+(45x+25)t .
\]
For $0\leq x\leq1$, use $t\geq1$ to get
$\Delta\geq x^2+5x+50>0$.
For $1\leq x<10/3$, use $t\geq x$ to get
$\Delta\geq46x^2-15x+25>0$.

Now suppose $t\leq\max(x,1)$. If $x\leq1$, then
$t\leq1$ contradicts $x>3-2t$. Hence $1<x<10/3$ and $t\leq x$.
The remaining demand after the larger endpoints is
$r=(5+4t-x)/2>0$.
If $r\leq3t$, then $H=tr$ and
\[
 \Delta=(x+5)^2+(45x-25)t>0.
\]
Otherwise
\[
 H=3t^2+x(5-x-2t)/2,\quad
 \Delta=26x^2-115x+25+(70x+100)t-50t^2 .
\]
Here $\partial_t\Delta=70x+100-100t\geq100-30x>0$.
By (R) it suffices to substitute
$t_0=\max((15-2x)/20,(3-x)/2)$.
For $1<x\leq15/8$ the resulting expression is
$(-43x^2+30x+125)/2$, a concave quadratic whose endpoint values
on $[1,15/8]$ are $56$ and $1925/128$.
For $15/8\leq x<10/3$ it is
\[
 \frac{148x^2-520x+575}{8}
 =\frac{37}{2}\left(x-\frac{65}{37}\right)^2+\frac{4375}{296}>0 .
\]
All cases have been covered.
\end{proof}

To prove part (2), apply Lemma~\ref{lem:alpha} if
$S\geq2N/5$ or $d+4c\geq2N/5$.
Otherwise $S>0$, so $a>0$.
Put $m=\min a_i$.
For $B_i$, with $u=S-a_i$, its mass and cost are
$b=u+4c$, $Q=cu$, and $b^2-16Q=(u-4c)^2$.
If any $B_i$ is heavy enough, scaling suffices.
In the remaining case,
\[
 d>S-2m-2c\geq3a-2c.
\]
Replace all leaves by their mean $a$, preserving $N$ and both
independent-set inequalities. Lemma~\ref{lem:star} gives a half on
the symmetric side $(d,a)+3(c,c)$.
Choose an actual largest leaf $M=\max a_i\geq a$.
The side $A_i$ for this leaf has capacities $(d,M)+3(c,c)$ and
contains that same selected vector: only one capacity has increased.
Its cost and mass are unchanged.
This proves the result without any further case split.
The case $S=0$ was already covered by $d+4c\geq2N/5$.

\section{A pentagon and a common nonneighbour}
Let $P=(v_0,\ldots,v_4)$ be an induced pentagon, let $z$ be nonadjacent
to all of $P$, and assign capacities $p_i,d,c$ as in part (3).
Put $W=\sum p_i$, $N=d+W+10c$, $h=N/2$.

\subsection{Geometry and direct cases}
Translate $z$ to $0$. Its ten nonneighbours correspond to the two-element
subsets $\{s,t\}$ of $D$, via the vertex $s+t$,
with adjacency given by disjointness: this is the Petersen
graph. For each $s\in D$, the four pairs containing $s$ are independent.
Their intersections with $P$ each have size at most two.
Each vertex of $P$ belongs to exactly two of these five four-sets,
so every intersection has size exactly two.
Consequently $z$ is in the three-point side of every pentagon matching cut.

For each $i$, that side $A_i$ consists of the four pairs
\[
 (p_i,p_{i+1}),\quad(p_{i+3},c),\quad(d,c),\quad(c,c).
\]
Set $s_i=p_{i+2}+p_{i+4}$.
Its complementary side has mass $s_i+6c$ and full cost $cs_i+2c^2$.
If it is heavy enough, scaling gives cost at most $N^2/64$.
Otherwise, for every nonadjacent pair of $P$,
\[
 s_i<(W+d-2c)/2. \tag{B}
\]
The independent set $N_H(z)$ has capacity $5c$.
Also one of the five independent four-sets above, together with $z$,
has capacity at least $d+2c+2W/5$.
Lemma~\ref{lem:alpha} therefore applies if $d+W\leq5c/2$ or
$d\geq10c/3$. This covers $c=0$.
In all remaining cases assume
\[
 c>0,\qquad 0\leq d\leq10c/3,\qquad W>5c/2-d. \tag{I}
\]

\subsection{Low centre capacity: $0\leq d\leq2c$}
Let $k=(W-d+2c)/2$.
Condition (B) implies every nonadjacent pair has mass less than $k$,
and every matching triple of $P$ has mass greater than $k$:
the omitted pair is less than $W-k=(W+d-2c)/2$.
In $A_i$, fill $z$, all four background classes, $p_i$, and $p_{i+3}$,
then take $k-p_i-p_{i+3}$ from $p_{i+1}$.
This is feasible and has mass $h$ and cost
\[
 q_i=p_i(k-p_i-p_{i+3})+cp_{i+3}+cd+c^2 .
\]
Cauchy--Schwarz gives the average bound
\[
 \frac15\sum_iq_i\leq H_{\rm low}
 =\frac{W^2}{50}-\frac{Wd}{10}+\frac{2cW}{5}+cd+c^2.
\]
The gap to the target is
\[
 G=\frac{N^2}{50}-H_{\rm low}
   =\frac{7Wd}{50}+\frac{d^2}{50}-\frac{3dc}{5}+c^2.
\]
The average pair mass is $2W/5$, so (B) implies $W>10c-5d$.
Set $r=d/c$, $w=W/c$.
Thus $w>\max(10-5r,5/2-r)$ and $0\leq r\leq2$.
The coefficient of $w$ in $G/c^2$ is nonnegative.
For $r\leq15/8$, substituting $w=10-5r$ gives
$1+4r/5-17r^2/25$, positive by concavity and endpoint values
$1,7/64$.
For $15/8\leq r\leq2$, substituting $w=5/2-r$ gives
$1-r/4-3r^2/25\geq1/50$.
Some $q_i$ therefore satisfies the target.

\subsection{Middle capacity and $W\geq c/2$}
Assume $2c\leq d\leq10c/3$.
In a suitable $A_i$, take $u=d/2+c\leq d$ from $z$, all four
background classes, and a half of the pentagon capacities.
The total is $h$ and the fixed cost is $cd/2+2c^2$.

If a nonadjacent pentagon pair has mass at least $W/2$, choose
mass $W/2$ within it. Each such pair lies in two matching triples,
with either member as the isolated vertex of the triple.
Explicitly, the pair $\{i,i+2\}$ lies in
$\{i+2,i+3,i\}$ and $\{i+4,i,i+2\}$, with isolated vertices
$i$ and $i+2$, respectively.
Make the smaller selected member the isolated vertex, so its mass
is at most $W/4$. The total cost is at most
\[
 H_{\rm pair}=cd/2+2c^2+cW/4 .
\]
A direct identity yields
\[
 \frac{N^2}{50}-H_{\rm pair}
 =\frac{(d-2c)^2}{50}
 +(W-c/2)\left(\frac W{50}+\frac d{25}+\frac{4c}{25}\right)\geq0 .
\]
If all nonadjacent pairs have mass less than $W/2$, the five explicit
constructions of Lemma~\ref{lem:c5} are feasible.
Their pentagon costs average at most $W^2/50$, and their background
attachment costs average $cW/5$. Some selection has cost at most
\[
 H_{\rm cycle}=W^2/50+cW/5+cd/2+2c^2.
\]
Here
\[
 \frac{N^2}{50}-H_{\rm cycle}
 =\frac{(d-2c)^2+c^2}{50}
 +(W-c/2)\left(\frac d{25}+\frac c5\right)>0 .
\]

\subsection{Middle capacity and $W\leq c/2$}
Retain (B) and (I).
In $A_i$ fill the three pentagon classes, $z$, and three background
classes. From the background partner of $z$ take
\[
 b_i=(4c-d-W)/2+s_i .
\]
It is positive because $d+W\leq23c/6<4c$, and is less than $c$ by (B).
The selection has mass $h$ and exact cost
\[
 F_i=p_ip_{i+1}+cp_{i+3}+c^2+
             d\big((4c-d-W)/2+s_i\big).
\]
Using $\sum s_i=2W$ and $\sum p_ip_{i+1}\leq W^2/2$, the average is
at most
\[
 H_{\rm small}=W^2/10+cW/5+c^2+d(4c-d-W)/2+2dW/5.
\]
With $r=d/c\in[2,10/3]$, $w=W/c\in[0,1/2]$, the normalized gap is
\[
 D=\frac{N^2/50-H_{\rm small}}{c^2}
  =\frac{13r^2}{25}+\frac{7rw}{50}-\frac{8r}{5}
   -\frac{2w^2}{25}+\frac w5+1.
\]
Its derivative in $w$ is at least $2/5$.
For $2\leq r\leq5/2$, use $w>5/2-r$ from (I) to get
$D\geq3r^2/10-21r/20+1\geq1/10$.
For $5/2\leq r\leq10/3$, use $w\geq0$ to get
$D\geq13r^2/25-8r/5+1\geq1/4$.
This proves part (3).

\section{An induced double star}
Use the representative in part (4), with centre capacities $d,e$,
left leaf capacities $a_1,a_2$, right leaf capacities $b_1,b_2$,
and common background capacity $c$.
Let $B(w)$ be the least cost of a half among the ten matching sides,
with infeasible sides assigned cost $+\infty$.

\begin{lemma}[Two matching-edge comparisons]\label{lem:smoothing}
For a fixed selection mass, averaging the two capacities of one matching
edge cannot decrease its minimum cost. For two disjoint edges
$(a,c),(b,c)$, averaging $a,b$ while keeping $c$ fixed also cannot
decrease the minimum cost.
These comparisons remain valid when other disjoint matching edges
are added.
\end{lemma}
\begin{proof}
For a single edge of total capacity $u$, put $L=\max(a,b)\geq u/2$.
At selected mass $t\leq u$, the cost is $L\max(0,t-L)$.
In its positive region its derivative in $L$ is $t-2L\leq0$.

For the second assertion fix $u=a+b$, write $a=u-b\geq b$, and increase
$b$ towards $u/2$. Let $g_b(t)$ be the minimum on these two edges.
If $a,b\leq c$, then $g_b(t)=c\max(0,t-2c)$, independent of $b$.
If $a,b\geq c$, put $q=\max(0,t-u)\leq2c$; then
\[
 g_b(t)=b\min(q,c)+(u-b)\max(0,q-c),
\]
whose coefficient of $b$ is nonnegative.
In the remaining region $b\leq c\leq a$, put $q=t-u-c\leq c$.
For $q\leq0$, the formula is $c\max(0,q+b)$.
For $q\geq0$, it is $cb+(u-b)q$, with coefficient $c-q\geq0$.
The formulas agree on their boundaries.
Thus $g_b(t)$ is nondecreasing throughout.
The compared component total capacities are unchanged. Taking the
infimum over mass allocations between this component and the remaining
disjoint edges preserves each comparison.
\end{proof}

Put $\ell=(a_1+a_2)/2$, $r=(b_1+b_2)/2$ and let $\bar w$ replace
the respective leaf pairs by $\ell,\ell$ and $r,r$.
Then
\[
 B(w)\leq B(\bar w),\qquad \alpha_w\geq\alpha_{\bar w}. \tag{S}
\]
For the matching comparison, the ten sides have four types:
the centre edge plus three background edges; all four leaf--background
edges; a centre--background edge, an opposite leaf--background edge,
and two background edges; or a centre--leaf edge, two opposite
leaf--background edges, and one background edge.
The first type is unchanged. For the second, apply
Lemma~\ref{lem:smoothing} to both leaf pairs.
For each side of the third type in $\bar w$, choose the larger actual
opposite leaf and use capacity inclusion.
For the fourth, first use the smoothing comparison on the two opposite
leaf--background edges, then choose the larger actual leaf on the centre
edge. Thus each side for $\bar w$ has a side for $w$ of no larger cost.
For the independence comparison, swapping bit coordinates $2,4$ fixes
both centres and the right leaves, and swapping coordinates $1,8$ fixes
both centres and the left leaves. They are automorphisms of $H$.
Averaging $w$ over these four automorphisms gives $\bar w$;
the maximum independent-set mass is a convex invariant function.

It suffices by (S) to treat the pair-balanced family.
The case $c=0$ is bipartite on its positive support, so an independent
half suffices. For $c>0$ scale to $c=1$, with
\[
 N=d+e+2\ell+2r+10,\qquad h=N/2 .
\]
The four-leaf matching side has mass $M=2\ell+2r+4$ and cost
$Q=2\ell+2r$, with $M^2-16Q=(2\ell+2r-4)^2$.
The two-exception side with centre $d$ and opposite leaf $r$ has
$M=d+r+6$, $Q=d+r+2$, and $M^2-16Q=(d+r-2)^2$.
The analogous side uses $e,\ell$.
If any is heavy enough, scaling proves the result.
Otherwise
\[
 d+e>2\ell+2r-2,\quad d-e<2\ell-2,\quad e-d<2r-2.
\]
Define positive parameters
\[
 k=(d-e)/2,\quad x=\ell-1-k,\quad y=r-1+k,\quad
 s=x+y,\quad z=(d+e)/2-1-s .
\]
Then $x,y,z>0$ and
\[
 \ell=1+x+k,\quad r=1+y-k,\quad
 d=1+s+z+k,\quad e=1+s+z-k,\quad h=8+2s+z.
\]
The independent sets $\{1,6,8,11,13\}$ and $\{2,4,7,9,14\}$
have masses $d+2r+2$ and $e+2\ell+2$.
Unless Lemma~\ref{lem:alpha} applies, with
\[
 \nu=x-y+k,\qquad H=(7-2s-z)/5=3-h/5
\]
their masses $h-3-\nu,h-3+\nu$ imply $|\nu|<H$ and $H>0$.

If $2s+z\leq2$, the centre-edge side has no greater minimum cost
than after averaging $d,e$ to $t=1+s+z$, by
Lemma~\ref{lem:smoothing}.
In $(t,t)+3(1,1)$, fill the large endpoints, the three background
small endpoints, and mass $1+s\leq t$ on the centre small endpoint.
Its cost is $3+(1+s)(1+s+z)$, and
\[
 N^2-50[3+(1+s)(1+s+z)]
 =56+28s+14z-34s(s+z)+4z^2
 \geq56-40s\geq16 .
\]
Here $s\leq1$ and $s+z\leq2$.

If $2s+z\geq2$, then $0<H\leq1$.
On the side meeting the double star in $\{1,2,6,8\}$, fill centre $d$,
the two opposite leaves $r,r$, and all four background classes.
Take $1+\nu=\ell-y$ on the remaining leaf of capacity $\ell$.
This quantity is positive and less than $\ell$, and the total is $h$.
Its cost is $F_L=d(1+\nu)+2r+1$.
The side meeting the double star in $\{1,2,4,9\}$ similarly gives
$F_R=e(1-\nu)+2\ell+1$, since $1-\nu=r-x$ is feasible.
Their average is
\[
 (F_L+F_R)/2=h-4+k\nu .
\]
As $k=\nu-x+y$, $|\nu|<H$, and $s\leq(7-5H)/2$,
\[
 k\nu\leq H^2+sH\leq H^2+\frac{H(7-5H)}2 .
\]
Since $h=15-5H$,
\[
 \frac{N^2}{50}-(h-4)=7-7H+2H^2 .
\]
Subtracting the preceding upper bound leaves
$\tfrac72(H-1)(H-2)\geq0$.
At least one of the two halves has cost at most $N^2/50$.
This completes the balanced family and hence part (4).
Lemma~\ref{lem:rounding} proves all integer assertions.

\section{A pentagon half avoiding a prescribed edge}
\begin{lemma}\label{lem:avoid}
For nonnegative pentagon capacities $b_0,\ldots,b_4$, with sum $T$,
there is a half of cost at most $T^2/50$ supported in
$\{i,i+1,i+3\}$ for some $i\ne0$.
Thus its template edge avoids the prescribed edge $\{0,1\}$.
\end{lemma}
\begin{proof}
If a nonadjacent pair has mass at least $T/2$, use it at zero cost.
Every such pair belongs to two pentagon templates, so at least one
avoids the prescribed edge.
Otherwise all the usual endpoint constructions are feasible.
Use the following four, with their indicated costs:
\[
\begin{array}{c|c}
\text{template}&\text{cost}\\
\{1,2,4\}&q_1=b_2(T/2-b_2-b_4)\\
\{4,0,2\}&q_2=b_4(T/2-b_4-b_2)\\
\{2,3,0\}&q_3=b_3(T/2-b_3-b_0)\\
\{3,4,1\}&q_4=b_3(T/2-b_3-b_1).
\end{array}
\]
Put $U=b_0+b_1$, $V=b_2+b_4$, $Z=b_3$. Direct expansion gives
\[
 \frac{T^2}{50}
 -\left(\frac{q_1+q_2}{5}+\frac{3(q_3+q_4)}{10}\right)
 =\frac{(2U-3V+2Z)^2+15(V-2Z)^2}{200}\geq0 .
\]
The coefficients are nonnegative and sum to one.
This proves the result, with the $T=0$ case immediate.
The endpoint constructions themselves are the known ones recalled in
Lemma~\ref{lem:c5}; the displayed identity provides the required
nonuniform averaging here.
\end{proof}

\section{A star with an extended leaf edge}
Give $0$ capacity $d$, the three simple leaves $1,2,4$ capacities
$a_1,a_2,a_3$, vertex $8$ capacity $e$, vertex $7$ capacity $t$,
and all other vertices capacity $c$. Put $a=(a_1+a_2+a_3)/3$.
Averaging these three leaves does not decrease the least cost among
the ten matching sides, and does not increase the maximum independent mass.
For independence, average over the permutations of bit coordinates
$1,2,4$, which fix $0,7,8$.
For matching, the sides with one simple leaf can choose the actual
largest leaf. A side with two simple leaves can omit the smallest;
the remaining pair has mean at least $a$, so first use
Lemma~\ref{lem:smoothing}, then capacity inclusion.
Sides with all three simple leaves use repeated equalizing transfers
between their leaf--background edges. Each transfer moves two capacities
towards their common mean without crossing, and has the monotonicity
proved in Lemma~\ref{lem:smoothing}.
Sides without simple leaves are unchanged.

It therefore suffices to treat $a_1=a_2=a_3=a$.
For $c=0$ the positive support is bipartite, so take an independent half.
For $c>0$ scale to $c=1$ and put
\[
 N=d+e+t+3a+10,\qquad h=N/2 .
\]
The six matching types are
\[
\begin{array}{ll}
 A=(d,a),(t,1),2(1,1),&B=2(a,1),(e,1),(1,1),\\
 C=(d,e),3(1,1),&D=3(a,1),(t,1),\\
 E=(d,1),3(1,1),&F=(e,t),3(a,1).
\end{array}
\]
Types $A,B$ each occur three times and are complementary;
$C,D$ and $E,F$ are complementary.
For $B,D,E$, their masses $M$ and costs $Q$ satisfy
\[
\begin{array}{c|c|c|c}
 &M&Q&M^2-16Q\\
 B&2a+e+5&2a+e+1&(2a+e-3)^2\\
 D&3a+t+4&3a+t&(3a+t-4)^2\\
 E&d+7&d+3&(d-1)^2 .
\end{array}
\]
Any heavy such side settles the bound by scaling.
Otherwise
\[
 2a+e+5<h,\quad 3a+t+4<h,\quad d+7<h. \tag{BM}
\]
Thus $A,C,F$ are heavy and $h>7$.
Four independent neighbourhoods have masses
\[
 d+t+3,\quad3a+e+1,\quad d+4,\quad e+4,
\]
witnessed respectively by $N_H(15),N_H(0),N_H(1),N_H(7)$.
If any reaches $2N/5$, use Lemma~\ref{lem:alpha}.
Otherwise, with
\[
 \nu=(d+t-3a-e+2)/2,\qquad H=3-h/5,
\]
the first two masses are $h-3\pm\nu$, so
$|\nu|<H$, $H>0$, $h<15$, and also $d,e<4h/5-4$.

For $7<h\leq10$, average the centre-edge capacities in $C$ to
$T=(d+e)/2$. This can only increase its minimum cost.
The resulting side $(T,T)+3(1,1)$ has capacity greater than $h$.
Its minimum cost is at most
\[
 U=3+(h-6)^2/4 .
\]
Indeed, for $T\geq1$ the larger-endpoint filling either costs at most $3$,
or costs $3+T(h-T-6)\leq U$.
For $T<1$, feasibility gives $T>(h-6)/2$ and $h<8$.
Fill both centre endpoints and then the background demand.
The cost is $h-4+(T-1)^2\leq U$.
Finally
\[
 N^2/50-U=(-17h^2+300h-1200)/100>0,
\]
a concave quadratic with endpoint values $67/100$ and $1$ on $[7,10]$.

For $10\leq h<15$, we have $0<H\leq1$.
On $A$, fill $d,t$ and the five background classes, then take
$1-\nu$ from the simple leaf. It is positive, and
$a-(1-\nu)=(d+t-a-e)/2>0$ by (BM).
On $F$, fill $e$, all three simple leaves and three background classes,
then take $1+\nu$ from the end vertex. It is positive, and
$t-(1+\nu)=(e+t+3a-d-4)/2>0$ by (BM).
Both have mass $h$. Their costs are
\[
 F_A=d(1-\nu)+t+2,\qquad F_F=e(1+\nu)+3a,
\]
with average $h-4+(e-d)\nu/2$.
The second inequality of (BM) gives $d+e>h-6$.
Combining this with $d,e<4h/5-4$ gives
$|e-d|<3h/5-2=7-3H$.
As $h=15-5H$, the target minus this average is at least
\[
 7-7H+2H^2-H(7-3H)/2
 =\tfrac72(H-1)(H-2)\geq0 .
\]
This proves part (5).

\section{A pentagon with a pendant vertex}
Let $P=(0,1,3,7,8)$ have capacities $p_0,\ldots,p_4$,
and give the pendant vertex $4$ capacity $d$.
The remaining ten capacities are $c$.
Put $W=\sum p_i$, $N=W+d+10c$, $h=N/2$.
Let $A_i$ be the matching side with pentagon triple
$\{i,i+1,i+3\}$, and $B_i$ its complement.
The pendant vertex belongs to $A_1,A_2,A_3$, and to $B_0,B_4$.
The pair types in $A_i$ are
\[
\begin{array}{ll}
 i=0,4:&(p_i,p_{i+1}),(p_{i+3},c),2(c,c),\\
 i=1,3:&(p_i,p_{i+1}),(p_{i+3},c),(d,c),(c,c),\\
 i=2:&(p_2,p_3),(p_0,d),2(c,c).
\end{array}
\]
Set $s_i=p_{i+2}+p_{i+4}$.
For $i=0,4$, $B_i$ has mass $s_i+d+5c$ and cost
$c(s_i+d)+c^2$, with square gap $(s_i+d-3c)^2$.
For the other three sides its mass is $s_i+6c$ and cost
$cs_i+2c^2$, with gap $(s_i-2c)^2$.
Unless a heavy $B_i$ already settles the bound, assume
\[
 s_0,s_4<(W-d)/2,\qquad
 s_1,s_2,s_3<(W+d-2c)/2. \tag{PB}
\]
All $A_i$ are then heavy.

First let $d\leq c$, and put
$k_0=(W+d)/2$, $k_1=(W-d+2c)/2$.
In $A_0,A_4$ fill the five background classes, and take pentagon mass $k_0$.
In $A_1,A_2,A_3$ fill $d$ and the four background classes, and take
pentagon mass $k_1$.
In each case fill $p_i,p_{i+3}$ and complete from $p_{i+1}$.
The omitted pair bounds in (PB) guarantee enough total capacity.
The filled nonadjacent pairs for $i=0,\ldots,4$ are
$s_1,s_2,s_3,s_4,s_0$; their bounds in (PB), using $d\leq c$,
guarantee nonnegative completion amounts.
The five costs are
\[
\begin{array}{ll}
 i=0,4:&q_i=p_i(k_0-p_i-p_{i+3})+cp_{i+3}+2c^2,\\
 i=1,3:&q_i=p_i(k_1-p_i-p_{i+3})+cp_{i+3}+cd+c^2,\\
 i=2:&q_2=p_2(k_1-p_2-p_0)+dp_0+2c^2 .
\end{array}
\]
Their sum equals
\[
 k_1W-\tfrac12\sum_i(p_i+p_{i+3})^2+cW
 +(d-c)(2p_0+p_4)+8c^2+2cd .
\]
Cauchy--Schwarz and $d-c\leq0$ give average at most
\[
 H_{\rm low}=W^2/50+2cW/5-Wd/10+(8c^2+2cd)/5 .
\]
The target gap is $7Wd/50+d^2/50+2c^2/5\geq0$.

Now let $d\geq c$.
Adding pairs of inequalities in (PB), using
$s_0+s_4=W-p_0$, $s_1+s_0=W-p_1$ and $s_3+s_4=W-p_4$,
gives $p_0>d$, $p_1>c$, $p_4>c$.
If $p_2\leq c$, the heavy side $A_1$ satisfies
\[
 Q=cM-4c^2+(p_1-c)(p_2-c)\leq cM-4c^2 .
\]
Scaling settles it. If $p_3\leq c$, use $A_3$ in the same way.
Hence only the case $p_i,d\geq c$ remains.

Write $b_i=p_i-c$, $q=d-c$, $T=\sum b_i+q$, so $N=16c+T$.
Use virtual pentagon capacities
$(b_0,b_1+q,b_2,b_3,b_4)$.
Lemma~\ref{lem:avoid} gives a half of cost at most $T^2/50$ in a
template with $i\ne0$.
Split its selected mass at vertex $1$ between the actual vertex $1$
and the pendant vertex $4$, with capacities $b_1,q$.
For $i=1,3$ the cycle vertex $0$ is omitted, so the pendant increment
has no additional increment edges, and lifting cannot increase cost.
For $i=2,4$, the selected mass at virtual vertex $1$ is zero.
Each lifted support lies in its actual matching side $A_i$.
Adding baseline $c$ at all eight side vertices gives mass $N/2$ and cost
at most $4c^2+cT/2+T^2/50\leq N^2/50$.
This proves part (6), including zero capacities and $c=0$.

\section{Two pentagons sharing a three-edge path}
Let the exceptional positions be $\{0,1,2,3,7,8\}$.
Write $d=w_0$, $e=w_3$, $t=w_7$, $f=w_8$.
Average $w_1,w_2$ to $a$.
A matching side containing both has them on separate background edges,
so Lemma~\ref{lem:smoothing} applies.
A side containing just one can use the actual larger capacity in its
interchanged counterpart. All other sides are unchanged.
Thus the least matching-side cost for the original weights is no greater
than for the averaged weights; the common target $N/2$ is unchanged.

For the balanced family $N=d+e+t+f+2a+10c$.
The relevant types are
\[
\begin{array}{ll}
 A=(d,a),(t,c),2(c,c),&A'=(a,e),(f,c),2(c,c),\\
 C=(e,t),(d,c),2(c,c),&D=2(a,c),(f,c),(c,c),\\
 E=(d,f),(e,c),2(c,c),&F=2(a,c),(t,c),(c,c),\\
 G=(d,c),(e,c),2(c,c),&H=(t,f),2(a,c),(c,c).
\end{array}
\]
Types $C,D$, $E,F$ and $G,H$ are complementary;
an $A$ also has a complementary $A'$.
Types $D,F,G$ satisfy $Q=cM-4c^2$.
Unless a heavy such side suffices, assume
\[
 2a+f+5c<h,\quad 2a+t+5c<h,\quad d+e+6c<h. \tag{TB}
\]
Adding the first and third inequalities, and then the second and third,
gives $t>c$, $f>c$.

If $e\leq c$, use the heavy side $C$:
its only non-background edge $(e,t)$ straddles $c$, so
$Q=cM-4c^2+(e-c)(t-c)\leq cM-4c^2$.
If $d\leq c$, similarly use $E$.
If $a\leq c$, one of the complementary sides $A,A'$ is heavy,
and both of their non-background edges straddle $c$.
All these cases give cost at most $N^2/64$.
We may therefore assume all $d,e,t,f,a\geq c$.

Use virtual pentagon capacities
\[
 b_0=d-c,\quad b_1=2(a-c),\quad b_2=e-c,\quad
 b_3=t-c,\quad b_4=f-c,
 \qquad T=\sum b_i,\quad N=16c+T .
\]
We require an extra property when lifting the virtual vertex $1$:
if it is isolated in the chosen template, its selected mass can be
split between the two actual parts;
if it is on the template edge, its selected mass must be at most
$b_1/2=a-c$, so one actual part can hold it.

If a nonadjacent pair has mass at least $T/2$, use a zero-cost half.
Pairs containing $1$ lie in the template $\{3,4,1\}$, where $1$ is isolated.
Pairs not containing $1$ lie in $\{2,3,0\}$ or $\{4,0,2\}$,
which omit $1$. Thus the extra property holds.

Otherwise all nonadjacent pair masses are below $T/2$.
If $b_2+b_4\leq b_0+b_3$, apply the four explicit constructions of
Lemma~\ref{lem:avoid} after shifting indices by one, thereby avoiding
edge $\{1,2\}$.
The only retained template with $1$ on an edge is $\{0,1,3\}$;
the chosen orientation fills $b_0,b_3$, so it takes
$T/2-b_0-b_3\leq b_1/2$ at $1$.
If $b_2+b_4\geq b_0+b_3$, use the unshifted lemma, avoiding $\{0,1\}$.
The only retained edge template using $1$ is $\{1,2,4\}$;
it fills $b_2,b_4$ and takes at most $b_1/2$ at $1$.
The square identity guarantees a half of cost at most $T^2/50$
among the respective four candidates.

Lift it as just described to the two actual parts at vertices $1,2$.
Templates $i=0,1$ use one such part, $i=2,4$ use none,
and $i=3$ allows both as isolated increment vertices.
Each support fits an actual matching side and its increment cost is
the virtual cycle cost.
Add baseline $c$ at all eight side vertices.
Again the mass is $N/2$ and the cost is at most
$4c^2+cT/2+T^2/50\leq N^2/50$.
This proves part (7) and completes Proposition~\ref{prop:families}.

\section{Proof of the six-exception theorem}
Enlarge the set of exceptional positions, if necessary, to a set $X$
of exactly six vertices; all other capacities are $c$.
If some generator colour is absent inside $X$, the missing-colour
argument in the five-exception proof gives cost at most $N^2/64$.

Otherwise select one edge of each of the five colours inside $X$.
Call the resulting graph $F$.
If it contains a cycle, its distinct edge colours sum to zero;
the unique nonempty zero-sum subset of $D$ forces a pentagon.
If $F$ is acyclic, its five edges on six vertices form a tree.
Any four-edge path uses four distinct colours; their sum is the fifth
generator, so its endpoints are adjacent in $H$, producing a pentagon.
A five-edge path is impossible, since its endpoint difference would be
the sum of all five generators, namely zero.

If $H[X]$ has no pentagon, this tree has diameter at most three.
It is therefore a star, an extended-leaf star, or a double star.
There are no additional edges in $H[X]$:
a putative new edge joins tree vertices at distance two or three;
the tree-path colours are distinct, and adding the new colour leaves
a nonempty set of at most four odd-multiplicity colours summing to zero,
a contradiction. Parts (2), (4), (5) apply.

If $H[X]$ contains a pentagon, its sixth vertex has at most two neighbours
on the cycle, since they must form an independent set on the pentagon.
Zero neighbours gives part (3), one gives part (6), and two nonadjacent
neighbours give part (7).

For completeness these abstract shapes have the representative labels
used above. Along any pentagon each of the five colours occurs once.
Every permutation of $D$ induces a linear automorphism: any four of its
generators are independent, and their unique relation is $\sum D=0$.
Together with translations this maps any cyclically ordered pentagon
to $(0,1,3,7,8)$.
The neighbours of $0$ outside this cycle are $2,4,15$:
the first also meets $3$, the last also meets $7$, and only $4$ is pendant.
For an extra vertex meeting $0,3$, the common neighbours of these
two vertices are exactly $1,2$, so the extra vertex is $2$.
The tree representatives follow by mapping their five distinct edge
colours and then translating a root.
The isolated-vertex proof already treats any common nonneighbour.

All cases therefore admit the required fractional half.
Lemma~\ref{lem:rounding} gives the integer assertion and proves
Theorem~\ref{thm:main}.


\section{A star and one further exceptional class}

\begin{theorem}\label{thm:sevenstar}
Let $z$ be a Clebsch vertex and let $t\notin\{z\}\cup N_H(z)$.
The seven capacities on $\{z\}\cup N_H(z)\cup\{t\}$ may be arbitrary
nonnegative real numbers. If the other nine capacities equal $c\geq0$,
there is a fractional half of cost at most $N^2/50$.
For integer capacities the corresponding blow-up has a set of
$\lfloor N/2\rfloor$ vertices inducing at most $N^2/50$ edges.
\end{theorem}
\begin{proof}
Translate $z$ to $0$. The ten nonneighbours of $0$ are the sums of two
distinct generators. Permuting the generators therefore sends $t$ to
$3=1+2$. Thus it suffices to use exceptional set
$X=\{0,1,2,3,4,8,15\}$, a four-cycle with three leaves at $0$.
For $c=0$ the supported graph is bipartite with parts $\{0,3\}$ and $D$;
a heavier part contains a zero-cost half. Hence normalize $c=1$.

First give vertices $0,3$ capacities $d,e$, vertices $1,2$ capacity $a$,
and vertices $4,8,15$ capacity $b$. All are nonnegative. Put
\[
 N=d+e+2a+3b+9=2h.
\]
The relevant matching types, with complementary pairs $A,B$ and $C,D$, are
\[
\begin{array}{ll}
 A=(d,a),3(1,1),&B=(a,e),3(b,1),\\
 C=(d,b),(e,1),2(1,1),&D=2(a,1),2(b,1).
\end{array}
\]
For a side of mass $M\geq h$ and full cost $Q$, proportional selection
costs $h^2Q/M^2$. In particular $Q\leq M^2/16$ gives $N^2/64$, while
$Q\leq2M^2/25$ gives $N^2/50$. Set
\[
 g(s)=2s^2/25-s+4
     =\frac2{25}(s-25/4)^2+\frac78\geq\frac78.
\]
The three actual independent sets
\[
 \{0,7,9,10,12\},\quad \{0,3,5,9,14\},\quad D
\]
have masses $d+4,d+e+3,2a+3b$. By Lemma~\ref{lem:alpha}, we may assume
\[
 d+4<4h/5,\qquad d+e+3<4h/5,\qquad2a+3b<4h/5. \tag{SI}
\]
Consequently $h>5+5d/4$ and $h<15$.
If $D$ is heavy, $Q_D=M_D-4$ and
$M_D^2-16Q_D=(M_D-8)^2$ settle the assertion. Otherwise
\[
 2a+2b+4<h,\qquad M_C=d+e+b+5>h. \tag{SD}
\]
For $C$, $Q_C=M_C-4+(d-1)(b-1)$.
If $(d-1)(b-1)\leq0$, scaling settles it by the same square identity.
Only the two same-side cases below remain.

\paragraph{Case I: $d,b\leq1$.}
On $C$, fill $d,b$, the background neighbour of $e$, and one endpoint
of each of the other two background edges. The initial mass $d+b+3$
is at most $5<h$. Complete from $e$ and the two unused background
endpoints. These additions all have unit cost per unit mass, and
heaviness ensures feasibility. The cost is
\[
 F_C=h-4+(1-d)(1-b). \tag{SC}
\]
If $h\geq15/2$, then
$g(h)-1=2(h-5)(h-15/2)/25\geq0$, so this suffices.
Assume $5<h<15/2$.

If $A$ is heavy and $a\geq1$, its exceptional edge straddles $1$;
$Q_A=M_A-4+(d-1)(a-1)\leq M_A-4$ suffices.
If $A$ is heavy and $a\leq1$, the same unit-completion construction gives
$F_A=h-4+(1-d)(1-a)$. With $u=\max(a,b)$, the better cost is
$h-4+(1-d)(1-u)$. From (SI), $2a+3b>6h/5-6$, hence
$u>6h/25-6/5$. Its target gap is at least
\[
 \frac{(2h-9)(h-5)+d(55-6h)}{25}\geq0.
\]
If instead $A$ is light, $d+a+6<h$ and the formula for $N$ imply
$e+3b>d+3$. Combining this with (SI) gives
$b>2+2d/3-4h/15$. The target gap in (SC) is therefore at least
\[
 \frac2{25}(h-15/2)(h-25/3)
 +d\left(\frac{4h}{15}-\frac13-\frac{2d}{3}\right)\geq0.
\]
The first product has two negative factors; the parenthesis is at
least $1/3$, since $h>5$ and $d\leq1$.

\paragraph{Case II: $d,b\geq1$ and $a\leq1$.}
If $A$ is heavy its edge $(d,a)$ straddles $1$, so scaling suffices.
Otherwise $B$ is heavy. If $e\geq1$, its edge $(a,e)$ straddles $1$
and again scaling suffices.
It remains that $e\leq1$ and $A$ is light.
Fill $a,e$ and all three $b$ endpoints in $B$.
By (SI) and $h>5$, their total mass satisfies
\[
 a+e+3b<4h/5-a+e\leq4h/5+1<h.
\]
Put $r=h-a-e-3b$. Heaviness of $B$ gives $0<r\leq3$.
Completing on its three unit endpoints gives cost
$F_B=ae+br$, and the identity
\[
 F_B-\bigl[h-4+(1-a)(1-e)\bigr]=(b-1)(r-3)\leq0
\]
shows $F_B\leq h-3$. This suffices if $h\geq15/2$.
If $h<15/2$, lightness of $A$ gives $d<3/2$, and
$2h=d+e+2a+3b+9<15$ with $d\geq1$ gives $b<5/3$.
Thus $(d-1)(b-1)<1/3<g(M_C)$, so
$Q_C\leq2M_C^2/25$ and proportional selection on $C$ suffices.

\paragraph{Case III: $d,b,a\geq1$ and $h\leq10$.}
Write $x=d-1$, $y=b-1$, and $A_0=a-1$, all nonnegative.
The total and (SD) give
\[
 15+e+x+2A_0+3y=2h\leq20,\qquad
 8+2A_0+2y<h\leq10.
\]
Hence $x+3y+e\leq5$ and $0\leq y<1$.
Here $M_C=7+e+x+y$ and $Q_C=M_C-4+xy$.
Since $g'(s)=4s/25-1>0$ for $s\geq7$,
\[
 \frac2{25}M_C^2-Q_C
 \geq\frac{f(x,y)}{25},\qquad
 f=2x^2-21xy+2y^2+3x+3y+23.
\]
For $0\leq y\leq3/4$, complete the square:
\[
 \frac{f(x,y)}{25}
 =\frac2{25}\left(x-\frac{21y-3}{4}\right)^2
  +\frac{7+6y-17y^2}{8}>0.
\]
The concave last term has endpoint values $7/8$ and $31/128$.
For $3/4\leq y\leq1$ and $0\leq x\leq5-3y$,
$f_x=4x-21y+3\leq23-33y<0$, so
\[
 f(x,y)\geq f(5-3y,y)=(1-y)(88-83y)\geq0.
\]
Thus proportional selection on the heavy side $C$ suffices.

\paragraph{Case IV: $d,b,a\geq1$ and $10\leq h<15$.}
Set $H=3-h/5$ and $\nu=d+e-h+6$.
By (SI), $|\nu|<H$ and $0<H\leq1$; by (SD), $b>1-\nu$.
On $C$, fill $d,e$ and all five background vertices, then take
$1-\nu$ from the $b$ endpoint. This lies in $[0,b]$ and gives mass $h$
and cost
\[
 F_C=h-4-(d-1)\nu. \tag{SN}
\]
For $\nu\geq0$ this is at most $h-4$.
Now put $u=-\nu\in(0,H)$.
If $e\geq1$, then $d-1\leq h-8-u$, so the error term in (SN)
is at most $(h-8-u)u$.
If $e\leq1$ and $B$ is heavy, its exceptional edge straddles $1$.
If $e\leq1$ and $B$ is light, then
$M_B=h+u+e-a<h$, so $a>e+u$.
On $A$, fill $d$ and all six background vertices and take $e+u$ from
the $a$ endpoint. This is feasible and has cost
\[
 F_A=h-4+(d-1)v,\qquad v=e+u-1\leq u<H,
 \qquad d-1=h-8-v.
\]
For $v\leq0$ this is at most $h-4$.
Otherwise both remaining errors have the form $(h-8-t)t$ with
$0\leq t\leq H$. This is increasing there because
$h-8-2H=7h/5-14\geq0$. Consequently the error is at most
\[
 (h-8-H)H=(7-6H)H.
\]
As $g(h)=7-7H+2H^2$, the target gap is at least
\[
 7-14H+8H^2=8(H-7/8)^2+7/8>0.
\]
This completes the balanced family.

\paragraph{Seven independent capacities.}
Now let $a_1,a_2$ be arbitrary capacities on $1,2$ and let
$b_1,b_2,b_3$ be arbitrary capacities on $4,8,15$.
Use means $a=(a_1+a_2)/2$ and $b=(b_1+b_2+b_3)/3$ only to select
one of the preceding constructions. The actual total $N$ is unchanged,
as are the masses of all three independent witnesses in (SI).
Every matching construction lifts as follows.

On $A$ choose a largest actual $a$ leaf and use only its mean capacity
$a$; on $C$ choose a largest actual $b$ leaf and use only its mean
capacity $b$. The other selected capacities are unchanged, so all the
corresponding proportional, unit-completion, and partial constructions
retain their mass and cost.
If balanced $D$ is heavy, omit a smallest actual $b$ leaf. Its actual
matching side has mass $M\geq2a+2b+4\geq h$ and full cost $M-4$;
the square identity gives a proportional half of cost at most $N^2/64$.

For $B$, choose a largest actual $a$ leaf, cap its usable capacity at
$a$, and take the matching side containing that leaf, vertex $3$, all
three actual $b$ leaves, and their three unit background partners.
Its mass $a+e+3b+3$ and full cost $ae+3b$ equal the balanced values.
Thus every proportional $B$ construction lifts unchanged.
For the explicit $B$ construction, fill $a,e$ and all three actual
$b$ leaves, then select $r/3$ from each background partner.
Because $0\leq r\leq3$, this is feasible and its cost is
$ae+(b_1+b_2+b_3)r/3=ae+br$, again exactly the balanced cost.

These sides exist in the full core: for generator $g$, the character
zero on $g$ and one on $D\setminus\{g\}$ has a zero fibre containing
$0,g$. For $g\in\{4,8,15\}$ it also contains $3=1+2$, giving $C$;
for $g\in\{1,2\}$ it omits $3$, giving $A$. Their complementary
fibres give $D$ and $B$, respectively. In the latter choose the
\emph{other} $a$ generator so that the largest $a$ leaf lies in $B$.
The required relabellings permute only the two $a$ leaves or the three
$b$ leaves and preserve all other exceptional roles.

Restoring $c$ multiplies masses by $c$ and costs by $c^2$.
The integer conclusion follows from Lemma~\ref{lem:rounding}.
\end{proof}

The theorem covers the 160 seven-position sets containing a closed
five-leaf star. The other seven-position shapes are not asserted here.
In particular this does not prove the arbitrary seven-exception theorem.

\section{A bounded-ratio criterion for all sixteen capacities}

\begin{proposition}\label{prop:ratio}
If all sixteen capacities are positive and
\[
 \frac{\max_vw_v}{\min_vw_v}\leq r_*:=\frac{24+5\sqrt{21}}{17},
\]
the Clebsch blow-up has a fractional half of cost at most $N^2/50$,
and the integer conclusion follows for integer capacities.
\end{proposition}
\begin{proof}
Write $m=\min_vw_v$, $M=\max_vw_v$, and $r=M/m$.
Consider any matching side, with pair totals $s_1,\ldots,s_4$,
total mass $L$ and full cost $Q$. Then
\[
 Q\leq\frac14\sum_js_j^2,\qquad
 \frac{Q}{L^2}\leq\frac{\sum_jt_j^2}{4(\sum_jt_j)^2},
 \quad t_j=s_j/(2m)\in[1,r].
\]
Holding the other coordinates fixed, and writing their sum and sum of
squares as $S,A$, gives derivative
\[
 \frac{d}{dt}\frac{t^2+A}{(t+S)^2}
 =\frac{2(tS-A)}{(t+S)^3}.
\]
An interior turning point is a minimum. Each coordinate can therefore
be moved to an endpoint without decreasing the ratio, so its maximum
over the box occurs at a corner. With $k$ coordinates equal to $r$,
the desired bound $Q/L^2\leq2/25$ follows from
\[
 P_k(r)=8(kr+4-k)^2-25(kr^2+4-k)\geq0.
\]
Here $P_0=28$, $P_4=28r^2$, and
\[
 P_1=-17r^2+48r-3,\quad
 P_2=P_1+(r-1)(15-r),\quad
 P_3=P_1+14(r^2-1).
\]
For $1\leq r\leq r_*<3$, all are nonnegative: $r_*$ is the larger
root of $P_1$ and its smaller root is below $1$.
Choose the heavier side of any generator cut. Its mass $L\geq N/2$,
so proportional selection gives cost at most
$(N/(2L))^2(2L^2/25)=N^2/50$.
Apply Lemma~\ref{lem:rounding} for integer capacities.
\end{proof}

The constant $r_*\approx2.75958108675$ is sharp for the particular
four-pair product inequality: one pair with both capacities $M$ and
three pairs with both capacities $m$ attain its limiting equality.
This is not a sharpness claim for the sparse-half problem itself.
The existence of some neighbourhood of the uniform weighting already
follows from its strict slack and continuity; no novelty is asserted
for that elementary observation or for the present explicit criterion.


\section{Verification and limitations}
The supplied exact-arithmetic verifier reconstructs the core, matching
cuts, and independent-set witnesses. It checks bounded examples of the
family representatives and verifies returned selections against all core edges.
Such finite checks support reproducibility; the proofs above establish
the unrestricted parameter claims.

The general sparse-half problem remains outside the scope of this note.
Uniform Clebsch graphs are covered by Razborov's strongly regular graph
result~\cite[Theorem 3.5 and Section 4.4]{R}; the present text concerns
specified nonuniform families. Known ingredients are attributed explicitly.
A bounded literature search has not located the exact six-exception or
seven-class star-extension statements, but this does not establish priority
or exclude a stronger prior theorem implying them.
The general seven-exception case and arbitrary weighted Clebsch graphs
remain unresolved here. There is no claim of a prize-qualifying solution.

The standard-library verifier \texttt{square\_star\_full\_check.py}
checks all construction branches on 12,472 bounded examples, returning
4,644 explicit fractional halves and performing 4,509 integer roundings.
The remaining 7,828 examples verify the premise of the attributed
independence-ratio theorem, rather than claiming to construct its half.
The separate \texttt{ratio\_check.py} checks exact corner-polynomial
identities and explicit bounded-ratio halves.

\paragraph{AI-use and contribution disclosure.}
This project was initiated and directed by the account
\texttt{kasaharamihano-spec}.
Proof development, literature retrieval, drafting, implementation, and
cross-checks were performed extensively with OpenAI Codex and multiple
AI agents. The listed account attribution must not be read as a claim that
the account holder personally derived or independently verified every step.
AI cross-checks are not external expert review or formal verification.
This draft is shared for examination with these limitations explicit.
No AI system is listed as a human author.
Exact symbolic workflow used SymPy 1.14.0 and the SymPy skill from
Scientific Agent Skills~\cite{K}; that is software-workflow attribution,
not a source for the mathematical claims.

\begin{thebibliography}{9}
\bibitem{B}
S. Brandt,
\emph{The local density of triangle-free graphs},
Freie Universit\"at Berlin, Preprint A 95-14, June 1995.
\href{https://ftp.mi.fu-berlin.de/pub/math/publ/pre/1995/index.html}{Official archive}.
\bibitem{NY}
S. Norin and L. Yepremyan,
\emph{Sparse halves in dense triangle-free graphs},
Journal of Combinatorial Theory, Series B \textbf{115} (2015), 1--25.
Author version: \href{https://arxiv.org/abs/1311.5818v2}{arXiv:1311.5818v2}.
\bibitem{R}
A. Razborov,
\emph{More about sparse halves in triangle-free graphs},
Sbornik: Mathematics \textbf{213} (2022), 109--128.
\href{https://doi.org/10.1070/SM9615}{doi:10.1070/SM9615};
author version: \href{https://arxiv.org/abs/2104.09406v2}{arXiv:2104.09406v2}.
\bibitem{K}
T. Kassis, V. Agarwal, Y. He, D. Patel and A. M. Brueckner,
\emph{Scientific Agent Skills: A Library of Procedural Knowledge for Research Agents},
2026. \href{https://doi.org/10.48550/arXiv.2609.00065}{arXiv:2609.00065}.
\end{thebibliography}
\end{document}
